\chapter{Introducción, problema y objetivos}\label{cap:introduccion}

\section{Motivación, contexto y problema}\label{sec:origen-motivacion}\label{sec:contexto-problema}\label{sec:arte-productos}\label{sec:arte-aportacion}

Goodreads y Letterboxd son dos servicios web de catalogación social: el primero, para libros,
y el segundo, para cine~\cite{goodreads,letterboxd}. Ambos comparten un patrón: una ficha canónica
por obra, independiente de la edición concreta que posea cada persona; un estado de consumo
personal, como pendiente, en curso o terminado; una valoración con una escala consistente;
listas ordenadas. Ni uno ni
otro representan la propiedad material con detalle, porque un libro o una película no plantean
el problema de inventario que sí plantea una consola con juegos en cartucho, disco y descarga.

En el ámbito del videojuego, Backloggd traslada casi literalmente el modelo de
Letterboxd~\cite{backloggd}, con buen tratamiento del backlog, pero sin inventario de
copias ni ambición de recomendación. OpenCritic agrega puntuaciones de prensa y es una
referencia de densidad visual y de tratamiento de la nota, más que un catalogador personal.
Steam ofrece una biblioteca personal rica, aunque acoplada a su tienda y limitada a lo
comprado en ella. La web de IGDB expone un catálogo amplio y bien modelado, orientado a la
consulta y no a la gestión de una colección. La Tabla~\ref{tab:productos} resume la comparación.

\begin{table}[htp]
\centering
\begin{tabularx}{\textwidth}{ | L{0.18\textwidth} | Z | Z | Z | Z | }
    \hline
    \textbf{Producto} & \textbf{Ficha canónica} & \textbf{Backlog} & \textbf{Inventario de copias} & \textbf{Recomendación} \\
    \hline
    Goodreads & Sí & Sí & No & Sí \\
    \hline
    Letterboxd & Sí & Sí & No & Parcial \\
    \hline
    Backloggd & Sí & Sí & No & No \\
    \hline
    OpenCritic & Sí & No & No & No \\
    \hline
    Steam & Parcial & Parcial & Solo digital de su tienda & Sí \\
    \hline
    IGDB.com & Sí & No & No & No \\
    \hline
    \textbf{SavePoint} & Sí & Sí & Sí & Sí \\
    \hline
\end{tabularx}
\caption{Comparación de SavePoint con los productos de referencia.}
\label{tab:productos}
\end{table}

El hueco que SavePoint ocupa se ve en la última columna de la
Tabla~\ref{tab:productos}: ningún producto reúne un inventario detallado de propiedad
material con un catálogo a escala real y un banco de pruebas de recomendación abierto y
reproducible. Es un problema propio del videojuego: un mismo título puede poseerse a la vez en
cartucho, en disco y en descarga, y en más de una plataforma, una distinción que no tiene ficha
ni equivalente en un libro o una película.

SavePoint se propone cubrir ese hueco. Aplica al videojuego el patrón de catalogación personal
de Goodreads y Letterboxd, con ficha canónica de cada obra, biblioteca, valoraciones y relaciones
de amistad, y añade lo que esos productos no necesitan y este medio sí: un inventario de copias
físicas y digitales por edición y plataforma, con procedencia de datos verificable en lugar de
un catálogo mantenido a mano. De ello se desprenden tres aportaciones. La primera es un
producto de catalogación de videojuegos que combina inventario de propiedad material, catálogo
a escala real y procedencia de datos documentada, un conjunto que ningún servicio existente
reúne, como muestra la Tabla~\ref{tab:productos}. La segunda es un banco de pruebas reproducible
para la comparación de recomendadores, con corpus citable, usuarios sintéticos y procedencia de
datos documentada, que se explota mediante un protocolo de evaluación congelado y una
comparación real entre dieciséis algoritmos, presentada en el Capítulo~\ref{cap:resultados}. La
tercera es la propia metodología de ingeniería asistida por agentes, aplicada de forma
sistemática y versionada en el repositorio, de manera que el proceso que produjo el sistema es
tan inspeccionable y reproducible por un tercero como el sistema mismo; el
Capítulo~\ref{cap:gestion} la desarrolla en profundidad.

El trabajo es, ante todo, un desarrollo de ingeniería del software: el diseño, la construcción,
la verificación y el despliegue de una plataforma web completa para catalogar videojuegos,
mantener una biblioteca personal y registrar valoraciones, listas, copias, comentarios y
relaciones sociales. El módulo de recomendación y su evaluación forman parte de esa
plataforma y tienen el mismo peso que el resto: sin un catálogo gobernado y una biblioteca con
señales reales no hay nada que recomendar, y sin una evaluación reproducible no hay forma
honesta de saber si lo recomendado funciona. La plataforma recoge y gobierna las señales que
los algoritmos consumen, y el laboratorio compara esos algoritmos sobre los mismos datos.

Una lista de recomendaciones no es evidencia suficiente de que un algoritmo funciona mejor que
otro. Sin un corpus congelado, una población definida, una división conocida y artefactos
identificables, dos ejecuciones pueden parecer comparables aunque en realidad respondan a
experimentos distintos, con datos distintos o con criterios de éxito distintos. SavePoint
separa por ello, de forma deliberada, la aplicación que presenta resultados a un usuario de los
trabajos offline que construyen y evalúan esos rankings. Esa separación no es solo una decisión
de arquitectura: se diseñó, implementó y validó mediante código, pruebas y evidencias
versionadas. Es esa separación la que permite que las conclusiones del
Capítulo~\ref{cap:resultados} de esta memoria se puedan volver a comprobar.

\section{Objetivos y alcance}\label{sec:objetivos}\label{sec:alcance-limites}

El objetivo general del trabajo es construir una plataforma de catalogación de videojuegos que
sirva, al mismo tiempo, como banco de pruebas reproducible para recomendación basada en
contenido, colaborativa e híbrida. De ahí se derivan las cuestiones que el trabajo
aborda: qué señales de un videojuego (género, plataforma, etiquetas curadas, calidad
externa, popularidad) se pueden representar de forma explicable en un catálogo real; cómo
cambia la recuperación de un ítem relevante cuando se combinan contenido, popularidad,
vecindad entre usuarios y diversidad de la lista resultante; y qué conclusiones son
defendibles cuando el corpus y la población de usuarios son, como se explica más abajo,
sintéticos y no personas reales.

Los objetivos específicos son gobernar el catálogo y su procedencia de datos; modelar
biblioteca, privacidad y relaciones sociales con el mismo nivel de cuidado que el catálogo;
implementar tanto los algoritmos de referencia (\textit{baselines}) como las tres familias de
recomendación descritas en el Capítulo~\ref{cap:estado-del-arte}; publicar un protocolo offline
con candidatas comunes para que la comparación entre algoritmos sea justa; y conservar hashes,
configuraciones y artefactos que hagan auditable cada resultado. La expectativa de partida del
trabajo no afirma la superioridad universal de ningún algoritmo: bajo el protocolo publicado,
combinar señales de contenido, popularidad y vecindad puede mejorar la recuperación frente a
los baselines más simples, pero esa mejora debe contrastarse sin confundir una diferencia
numérica con significación estadística o con validez fuera del corpus y la población con los
que se midió.

El alcance funcional del producto incluye catálogo, importación gobernada de datos, biblioteca,
copias físicas y digitales, valoraciones, listas, comentarios, perfiles, privacidad, módulo
social, recomendaciones y las medidas de seguridad y copia de
seguridad propias de un servicio con usuarios. De las dieciséis variantes de recomendación
evaluadas, la aplicación publica tres, según se explica en la
Sección~\ref{sec:alg-publicadas}. El alcance de la evaluación, en cambio, se limita a
una publicación de evidencia concreta, identificada como v15 y descrita en el
Capítulo~\ref{cap:resultados}, con una única ejecución del test. Conviene subrayar que estos usuarios sintéticos no son personas
reales, sino perfiles generados de forma controlada para poder ejercitar el sistema de
recomendación sin depender de una base de usuarios real que este TFG no tiene tiempo ni alcance
para reunir. Los resultados, por tanto, dependen del corpus concreto, de la semilla aleatoria,
de la división de datos y del protocolo empleados: no permiten generalizar automáticamente a
preferencias de jugadores reales ni declarar un método como mejor fuera de esas condiciones.

\section{Estructura de la memoria}\label{sec:estructura-memoria}

La memoria sigue el recorrido natural de un proyecto de ingeniería del software, tratando la
plataforma y el estudio de recomendadores en un único hilo. El
Capítulo~\ref{cap:estado-del-arte} revisa las fuentes de datos, los sistemas de
recomendación y su evaluación, y el método de trabajo. El Capítulo~\ref{cap:gestion} describe la
gestión, la metodología y la planificación del trabajo, con su calendario, su presupuesto y sus
riesgos. El Capítulo~\ref{cap:analisis} recoge el análisis del problema: actores, requisitos de
la plataforma y de la evaluación de recomendadores, reglas de negocio y trazabilidad. El
Capítulo~\ref{cap:diseno} presenta el diseño y la arquitectura, incluidos los algoritmos de
recomendación con sus fórmulas y ponderaciones. El Capítulo~\ref{cap:implementacion} describe la
implementación, los datos, las pruebas y el despliegue. El Capítulo~\ref{cap:resultados}
somete las dieciséis variantes al protocolo de evaluación y presenta los resultados. El
Capítulo~\ref{cap:conclusiones} cierra con el grado de cumplimiento de los objetivos, las
limitaciones y las líneas de continuación.
